%% LyX 2.0.8.1 created this file.  For more info, see http://www.lyx.org/.
%% Do not edit unless you really know what you are doing.
\documentclass[english,aps, print]{revtex4-1}
\usepackage{charter}
\usepackage[T1]{fontenc}
\usepackage[latin9]{inputenc}
\setcounter{secnumdepth}{3}
\usepackage{color}
\usepackage{babel}
\usepackage{amsmath}
\usepackage{amssymb}
\usepackage{graphicx}
\usepackage{esint}
\usepackage[unicode=true,
 bookmarks=true,bookmarksnumbered=false,bookmarksopen=false,
 breaklinks=false,pdfborder={0 0 0},backref=false,colorlinks=true]
 {hyperref}
\hypersetup{pdftitle={Validation of Flaps2D simulation code},
 pdfauthor={Thibault J.-Y. Derrien <derrien@fzu.cz>},
 citecolor=blue, urlcolor=blue, linkcolor=blue}

\makeatletter

%%%%%%%%%%%%%%%%%%%%%%%%%%%%%% LyX specific LaTeX commands.
\newcommand{\noun}[1]{\textsc{#1}}
%% Because html converters don't know tabularnewline
\providecommand{\tabularnewline}{\\}

\makeatother

\begin{document}

\title{Validation of the Flaps 2D simulation code}


\author{Thibault J.-Y. Derrien}

\email{derrien@fzu.cz}


\affiliation{HiLASE Centre, Institute of Physics, Academy of Science of the Czech
Republic, \\
Za Radnic� 828/5, 25241 Doln� B\v{r}e�any, Czech Republic}


\date{August 30th, 2015}
\begin{abstract}
This report aims to document and verify the ``Flaps 2D'' simulation
tool, developed by myself since February 2012, to investigate the
physics of laser-matter interaction in nanostructured materials. 
\end{abstract}
\maketitle

\section*{Context of the simulation tool}

The development of this work has been supported during 11 months by
the french National Research Agency (ANR) under the project ``Ultrasonde''
leaded by Tatiana Itina, Laboratoire Hubert Curien, Saint-Etienne,
France. 

Today, I'm still working on the development of this code which, at
term, allows to study the absorption by a structured material irradiated
by a femtosecond laser. The code has been designed so that material
can be also polarized by a high voltage. The coupling of polarized
materials with laser illumination allow tomographic reconstruction
of atomic lattices in 3D with metals, dielectrics and oxides. The
code can also be used to study the transport of free-carriers in nanostructures
considering the effect of the geometry. This aspect is also a promiseful
domain for a better understanding of the carrier and heat transport
properties at the nanoscale. 

\tableofcontents{}

\clearpage{}


\section{System of equations solved by the code}

The drift induced by the external field and by the charge separation
is neglected. 


\paragraph{Free-carrier generation and transport}

{\small{}
\begin{equation}
\begin{cases}
\frac{\partial n_{e}}{\partial t}+\boldsymbol{\nabla}\cdot\boldsymbol{J_{e}} & =G_{e}-R_{e}\\
\frac{\partial n_{h}}{\partial t}+\boldsymbol{\nabla}\cdot\boldsymbol{J_{h}} & =G_{h}-R_{h}.
\end{cases}\label{eq:Transport}
\end{equation}
}
\[
G_{e,h}=\sigma_{1}I+\sigma_{2}I^{2}+\delta_{II}^{e,h}n_{e,h}
\]
\[
R_{e,h}=R_{h}^{A}+R_{e}^{A}=C_{h}^{AR}n_{h}^{2}n_{e}+C_{e}^{AR}n_{e}^{2}n_{h}
\]
{\small{}
\[
\boldsymbol{J_{e,h}}=-D_{e,h}\boldsymbol{\nabla}n_{e,h}
\]
\[
D_{e,h}=\mu_{e,h}\frac{k_{B}T_{e,h}}{e}\frac{\mathcal{F}_{1/2}\left(\eta_{e,h}\right)}{\mathcal{F}_{-1/2}\left(\eta_{e,h}\right)},
\]
 
\[
\mu_{e,h}=\frac{q_{e,h}}{m_{e,h}^{\star}\nu_{e,h}}\frac{\mathcal{F}_{0}\left(\eta_{e,h}\right)}{\mathcal{F}_{1/2}\left(\eta_{e,h}\right)}
\]
}{\small \par}


\paragraph{Free-carrier energy }

{\small{}
\begin{equation}
\begin{cases}
\frac{\partial}{\partial t}\left(C_{e}T_{e}\right) & +\boldsymbol{\nabla}\left(\kappa_{e}\boldsymbol{\nabla}T_{e}\right)=Q_{e}-\gamma_{e-ph}(T_{e}-T_{s})\\
\frac{\partial}{\partial t}\left(C_{h}T_{h}\right) & +\boldsymbol{\nabla}\left(\kappa_{h}\boldsymbol{\nabla}T_{h}\right)=Q_{h}-\gamma_{h-ph}(T_{h}-T_{s})
\end{cases}\label{eq:TTMe}
\end{equation}
}{\small \par}

{\small{}
\[
\frac{\partial\left(C_{Si}T_{Si}\right)}{\partial t}=\nabla\left(\kappa_{Si}\nabla T_{Si}\right)+\gamma_{e-ph}(T_{e}-T_{Si})+\gamma_{h-ph}(T_{h}-T_{Si})
\]
}{\small \par}

{\small{}
\[
C_{e,h}=\frac{3}{2}k_{B}n_{e,h}\left\{ \frac{\mathcal{F}_{3/2}\left(\eta_{e,h}\right)}{\mathcal{F}_{1/2}\left(\eta_{e,h}\right)}-\eta_{e,h}\left[1-\frac{\mathcal{F}_{3/2}\left(\eta_{e,h}\right)}{\mathcal{F}_{1/2}\left(\eta_{e,h}\right)}.\frac{\mathcal{F}_{-1/2}\left(\eta_{e,h}\right)}{\mathcal{F}_{1/2}\left(\eta_{e,h}\right)}\right]\right\} 
\]
}{\small \par}

{\small{}
\[
\kappa_{e,h}=\frac{k_{B}^{2}n_{e,h}\mu_{e,h}T_{e,h}}{e}\left[6\frac{\mathcal{F}_{2}\left(\eta_{e,h}\right)}{\mathcal{F}_{0}\left(\eta_{e,h}\right)}-4\frac{\mathcal{F}_{1}\left(\eta_{e,h}\right)^{2}}{\mathcal{F}_{0}\left(\eta_{e,h}\right)^{2}}\right]
\]
\begin{align*}
Q_{e,h} & =\frac{\hbar\omega-E_{gap}}{\hbar\omega}\sigma_{1}I+\frac{2\hbar\omega-E_{gap}}{2\hbar\omega}\sigma_{2}I^{2}\\
 & -E_{gap}\delta_{II}(T_{e,h})n_{e,h}+\alpha_{fcr}^{e,h}I+E_{gap}R_{e,h}^{A}
\end{align*}
}{\small \par}


\paragraph{Lattice energy}

\begin{equation}
\frac{\partial\left(C_{s}T_{s}\right)}{\partial t}-\boldsymbol{\nabla}\left(\kappa_{s}\boldsymbol{\nabla}T_{s}\right)=\gamma_{e-ph}(T_{e}-T_{s})+\gamma_{h-ph}(T_{h}-T_{s})\label{eq:TTMlattice}
\end{equation}


\begin{equation}
\gamma_{e,h-ph}=C_{e,h}\nu_{e,h-ph}.\label{eq:Coupling}
\end{equation}
\[
\nu_{e,h-ph}^{-1}=\tau_{ph}^{0}\left[1+\left(\frac{n_{e,h}}{n_{th}}\right)^{2}\right]
\]


\[
C_{s}\left(T_{s}\right)=10^{3}\rho_{Si}\times\left(a_{0}T_{s}+a_{1}T_{s}^{a_{2}}\right)
\]
\[
\kappa_{s}(T_{s})=10^{2}\left(a_{0}+\frac{a_{1}}{1+e^{-\frac{T_{s}-T_{c}+\frac{w_{1}}{2}}{w_{2}}}}\right)\left(1-\frac{1}{1+e^{-\frac{T_{s}-T_{c}-\frac{w_{1}}{2}}{w_{3}}}}\right)
\]


Since the Fermi-Dirac statistics strongly reduce the possible comparison
with analytical solutions, it will be neglected in the following Sections
\ref{sec:Free-carrier-density:-Validation}, \ref{sec:Free-carrier-energy:-validation}
and \ref{sec:Coupled-effects}. However, it will be considered in
the section \ref{sec:Control-of-the}. 


\section{Free-carrier density: Validation of analytical cases\label{sec:Free-carrier-density:-Validation}}


\subsection{Auger recombination}

The formulation proposed by \cite{Silaeva2012a} has to be validated
within several steps. Auger recombination is modeled by a coupled
non-linear system. In that section, we only consider the equation
of free-carrier given by

\begin{equation}
\frac{\partial n_{e}}{\partial t}=\frac{\partial n_{h}}{\partial t}=C_{h}^{AR}n_{h}^{2}n_{e}+C_{e}^{AR}n_{e}^{2}n_{h}\label{eq:AugerNe}
\end{equation}
Analytical solution of the equation is given by 
\[
n_{e}(t)=\frac{C_{AR}^{(h)}n_{e}^{0}n_{h}^{0}}{\left(n_{e}^{0}C_{AR}^{(e)}+n_{h}^{0}C_{AR}^{(h)}\right)\exp\left(C_{AR}^{(h)}n_{h}^{2}t\right)-C_{AR}^{(e)}n_{e}^{0}}
\]


\begin{figure}[h]
\begin{centering}
\includegraphics[width=8cm,bb = 0 0 200 100, draft, type=eps]{home/thibault/Documents/LaAPT/Flaps2D/09-EnergyConservation/20140314-ValidationOfAuger.eps}
\par\end{centering}

\caption{\label{fig:AugerNe}Relative error calculation (\%) of Eq. (\ref{eq:AugerNe})
as a function of time. }
\end{figure}
Figure \ref{fig:AugerNe} presents the comparison between the numerical
calculation and the analytical solution for an initial value problem,
where densities have been fixed at $N_{e,h}(t=0)=10^{27}$ m$^{-3}$.
The mismatch between numerical and analytical calculations is smaller
than $0.005$\%. 


\subsection{Interband absorption}


\paragraph{One-photon absorption}

Validation on analytical function at $\lambda=343$ nm has been performed
but is not yet presented. 


\paragraph{Two-photon absorption}

Test at $\lambda=1030$ nm. Test is shitty ! Go to second order time
integration scheme. 


\subsection{Impact ionization}

To be completed. 


\section{Free-carrier energy: validation of analytical cases\label{sec:Free-carrier-energy:-validation}}


\subsection{Heating of the interband absorption}

In this subsection, we simply consider the interband absorption at
343 nm. 
\[
\frac{\partial}{\partial t}\left(C_{e}T_{e}\right)=\frac{\hbar\omega-E_{gap}}{\hbar\omega}\sigma_{1}I
\]
where the free-carrier density has been fixed to $n_{e}=10^{27}$
m$^{-3}$. 

\begin{figure}[h]
\begin{centering}
\includegraphics[width=8cm,bb = 0 0 200 100, draft, type=eps]{home/thibault/Documents/LaAPT/Flaps2D/09-EnergyConservation/20130905-HeatingIonization-1E27.eps}
\par\end{centering}

\caption{\label{fig:ErrorSigma1-Heating}Relative error between numerical and
analytical solutions for the interband absorption of electrons. $\lambda=343$
nm. }
\end{figure}
Figure \ref{fig:ErrorSigma1-Heating} shows the relative error between
the numerical and the analytical solution. The maximum relative error
is $0.1$\%, which is reasonable. This relative error can be reduced
by replacing the first-order upwind time integrator by a second-order
one. 

At 1030 nm, the two-photon absorption plays an important role. It
is thus of interest to validate the integrator at this wavelength.
For that purpose, we push the laser fluence to $100$ J.m$^{-2}$
in order to reach a good 2-photon absorption probability. Figure \ref{fig:TwoPhotonTeError}
presents the free-carrier heating by the 2-photon absorption process
at $n_{e}=10^{27}$ m$^{-3}$. 
\begin{figure}[h]
\begin{centering}
\includegraphics[width=8cm,bb = 0 0 200 100, draft, type=eps]{home/thibault/Documents/LaAPT/Flaps2D/09-EnergyConservation/20130905-HeatingIonization2-1E27.eps}
\par\end{centering}

\caption{\label{fig:TwoPhotonTeError}Relative error between numerical and
analytical solutions for the calculation of 2-photon absorption at
$\lambda=1030$ nm wavelength, $F=100$ J.m$^{-2}$ and pulse duration
$\tau=40$ fs. }


\end{figure}



\subsection{Validation of the electron energy coupling with lattice}


\paragraph{Cooling of the electrons by the cold lattice}

In this section, the equations are reduced to the following{\small{}
\begin{equation}
\frac{\partial}{\partial t}\left(C_{e}T_{e}\right)=-\gamma_{e-ph}(T_{e}-T_{s})\label{eq:TTMe-2}
\end{equation}
}{\small \par}

\begin{equation}
\gamma_{e,h-ph}=C_{e,h}\nu_{e,h-ph}.\label{eq:Coupling-1}
\end{equation}
\[
\nu_{e,h-ph}^{-1}=\tau_{ph}^{0}\left[1+\left(\frac{n_{e,h}}{n_{th}}\right)^{2}\right]
\]


Free-carrier density $n_{e}$ is taken as a constant $n_{e}=10^{27}$
m$^{-3}$. The initial carrier temperature is set to the value of
$T_{e}=300,000$ K. 

\begin{figure}[h]
\begin{centering}
\includegraphics[width=8cm,bb = 0 0 200 100, draft, type=eps]{home/thibault/Documents/LaAPT/Flaps2D/09-EnergyConservation/20140317-CheckCoupling.eps}
\par\end{centering}

\caption{\label{fig:Relative-error-on-Te-Coupling}Relative error on the calculation
of the free-carrier temperature. In this result, only the thermal
coupling between electrons and lattice has been checked. The sudden
change of the relative error is due to the change of timestep during
the calculation. }
\end{figure}
Figure \ref{fig:Relative-error-on-Te-Coupling} presents the relative
error between the numerical and the analytical solution as a function
of time. Results show that the free-carrier temperature is calculated
with an error lower than $0.01$ \% from the exact value. The effect
of the change of time-step is also observed: the change at $t=10^{-12}$
s induces a less precise calculation, which is then amplified until
the end of the interaction. The overall error remains however lower
than $0.015$ \%, which is rather good. 


\section{Validation of lattice thermal coupling}


\subsection{Heating of the lattice by hot electrons}


\paragraph{Reduction to a ``simple'' case}

This section investigates the heating of the lattice in our conditions.
Eq. \ref{eq:TTMlattice} is reduced to its simplest form allowing
to validate to calculate the heating of the lattice. The laser source
is disabled, and we only consider the coupling between a gas of hot
electrons at a fixed density coupled with the lattice. This allows
to set the carrier density $n_{e}$ as a constant, and the free-carrier
temperature $T_{e}$ as a constant. Thus, the carrier heat capacity
$C_{e}$, and the electron-lattice coupling rate $\gamma_{e-ph}$
are constant. These approximations lead to set the following equation 

\begin{equation}
\frac{\partial\left[C_{s}(T_{s}(t))\times T_{s}(t)\right]}{\partial t}=\gamma_{e-ph}\left[T_{e}^{0}-T_{s}(t)\right]\label{eq:TTMlattice-1}
\end{equation}
where 
\begin{equation}
\gamma_{e,h-ph}=C_{e,h}\nu_{e,h-ph}.\label{eq:Coupling-1-1}
\end{equation}
\[
\nu_{e,ph}^{-1}=\tau_{ph}^{0}\left[1+\left(\frac{n_{e}^{0}}{n_{th}}\right)^{2}\right].
\]
In our case, the lattice heat capacity strongly depends to the lattice
temperature $T_{s}$. Depending on the form of the coefficient $C_{s}\left(T_{s}\right)$,
Eq. \ref{eq:TTMlattice-1} will therefore exhibit strong non-linearities.
Before considering solving this equation, we first develop the possible
functions used to adjust the lattice heat capacity on the available
experimental results. In all the cases, the equation is non-linear
and does not admit a simple analytical solution. We are thus able
to compare the quality of the numerical integration with other numerical
integration packages and keep checking the energy conservation of
our results. However, no error estimation will be possible to perform
since no analytical solution is available. 


\paragraph{Adjustment of the lattice heat capacity}

Here is a list of the functions which have been built as a function
of lattice temperature fitting the experimental measurements of Refs.
\cite{Flubacher1959,Okhotin1972}. Each case will lead to different
types of non-linearities. Since we have different possibilities to
adjust the experimental data, the type of function should be chosen
in order to minimizer the difficulties over the numerical integration. 
\begin{enumerate}
\item Exponent fraction
\[
C_{s}(T_{s})=\frac{a_{0}}{e^{a_{1}/T_{s}}-a_{2}}=\frac{0.4129\times10^{3}\rho_{Si}}{e^{\frac{88.183}{T_{s}}}-0.6765}
\]
The resulting type of the equation is 
\[
a_{0}a_{1}\left(\frac{\partial{\it T_{s}}}{\partial t}\right){\rm e}^{\frac{a_{1}}{{\it T_{s}}\left(t\right)}}\left({\rm e}^{\frac{a_{1}}{{\it T_{s}}\left(t\right)}}-a_{2}\right)^{-2}\left({\it T_{s}}\left(t\right)\right)^{-1}+a_{0}\frac{\partial T_{s}}{\partial t}\left({\rm e}^{\frac{a_{1}}{T_{s}\left(t\right)}}-a_{2}\right)^{-1}=\frac{{\it C_{e}}\left({\it T_{e}^{0}}-{\it T_{s}}\left(t\right)\right)}{{\it \tau_{e}}}
\]

\item Sum of exponents
\[
C_{s}(T_{s})=\eta\left(ae^{bT_{s}}+ce^{dT_{s}}\right)
\]
where $\eta=10^{3}\rho_{Si}$, $a=0.899$, $b=5.455\times10^{-5}$,
$c=-0.959$, and $d=-0.004218$. The resulting type of equation is
\[
\left(\frac{\partial{\it T_{s}}}{\partial t}\right)\eta\left({\it T_{s}}\left(t\right)abe^{b\times{\it T_{s}}\left(t\right)}+{\it T_{s}}\left(t\right)cd{e}^{d\times{\it T_{s}}\left(t\right)}+ae^{b\times{\it T_{s}}\left(t\right)}+ce^{d\times{\it T_{s}}\left(t\right)}\right)=\frac{{\it C_{e}}\left(-{\it T_{e}^{0}}+{\it T_{s}}\left(t\right)\right)}{\tau_{e}}
\]

\item Real-valued polynomial
\[
C_{s}\left(T_{s}\right)=\sum_{i=0}^{i=5}a_{i}T_{s}^{i}
\]
 where $a_{5}$ = $2.36\times10^{-16}$, $a_{4}=-1.707\times10^{-12}$,
$a_{3}=4.619\times10^{-9}$, $a_{2}=-5.912\times10^{-6}$, $a_{1}=3.733\times10^{-3}$,
and $a_{0}=-0.0494$. The resulting of equation is 
\[
\left(\frac{\partial T_{s}}{\partial t}\right)\left(2a_{1}{\it T_{s}}\left(t\right)+3a_{2}\left({\it T_{s}}\left(t\right)\right)^{2}+4a_{3}\left({\it T_{s}}\left(t\right)\right)^{3}+5a_{4}\left({\it T_{s}}\left(t\right)\right)^{4}+6a_{5}\left({\it T_{s}}\left(t\right)\right)^{5}+a_{0}\right)=\frac{{\it C_{e}}\left(-{\it T_{e}^{0}}+{\it T_{s}}\left(t\right)\right)}{\tau_{e}}
\]

\item Van Driel's style fit function
\[
C_{s}\left(T_{s}\right)=10^{3}\rho_{Si}\times\left(a_{0}T_{s}+a_{1}T_{s}^{a_{2}}\right)
\]
with $a_{0}=-0.003592$, $a_{1}=0.01458$, and $a_{2}=0.8316$. The
stability of this function is very good. However, the temperature
validity range of this function is limited to 1400 K. 
\end{enumerate}
\begin{figure}[h]
\centering{}\includegraphics[width=12cm]{\string"/media/thibault/PHD/Travail/Post Doc/Calculs/CiKi-Silicon/20131101-CiModels\string".pdf}\caption{Heat capacity of Si and various adjustments. }
\end{figure}



\paragraph{Numerical integration}

Several possibilities can be tried to deal with the non-linearity. 
\begin{enumerate}
\item Use a higher-order precision scheme on the time integration for the
$T_{s}$ equation. 
\item Use a Runge-Kutta solver. However, there gonna be diffusive terms
enabled later, which will not be mathematically clean with such a
scheme. 
\item Iterate the solution and ensure convergence \emph{before going to
the next time-step }(can really multiply the calculation time)\emph{.}
\end{enumerate}
Applying a higher-order precision scheme allows to obtain solutions
closer to the ``exact'' solution. Here, the ``exact'' solution
is considered to be the one given by the 4th order Runge Kutta scheme
(``ode45'' function of Maple and Matlab). However, using an eulerian
scheme does not provide a sufficiently good convergence of the solution
even at the 4th order. The choice of the fitting function using the
Van Driel's form interestingly allows to obtain a solution equal to
the ``exact'' one. The origin of this numerical recipe is not understood
yet. 

It should be mentioned that in the all the cases considered in this
section, the energy conservation has been verified and did not show
any deficiency. The serious non-linearity that we face in this section
is important for the topic and can strongly affect the final results
given by the code. 


\section{Verification of the coupled effects\label{sec:Coupled-effects}}


\subsection{Change of carrier heat capacity}


\subsection{Auger recombination}

\[
\frac{\partial n_{e}}{\partial t}=\frac{\partial n_{h}}{\partial t}=C_{h}^{AR}n_{h}^{2}n_{e}+C_{e}^{AR}n_{e}^{2}n_{h}
\]


\begin{align}
\frac{\partial}{\partial t}\left(C_{e}T_{e}\right) & =E_{gap}C_{AR}^{(e)}N_{e}^{2}N_{h}\\
\frac{\partial}{\partial t}\left(C_{h}T_{h}\right) & =E_{gap}C_{AR}^{(h)}N_{h}^{2}N_{e}
\end{align}
If $n_{e}(t=0)=n_{h}(t=0)$, then $n_{e}(t)=n_{h}(t)$. The system
is thus equivalent to the following 
\begin{equation}
\begin{cases}
\frac{\partial n_{e}}{\partial t} & =C_{AR}^{(e)}n_{e}^{3}\\
\frac{\partial}{\partial t}\left(C_{e}T_{e}\right) & =E_{gap}C_{AR}^{(e)}N_{e}^{3}
\end{cases}\label{eq:AugerCoupled}
\end{equation}
with the initial conditions $T_{e,h}(t=0)=T_{e,h}^{(0)}$ and $N_{e}(t=0)=N_{e}^{(0)}$.
The solution of the Eq. (\ref{eq:AugerCoupled}) is 
\[
\begin{cases}
N_{e}(t) & =\left(\frac{1}{N_{e,0}^{2}}+2C_{AR}^{(e)}t+2C_{AR}^{(h)}t\right)^{-\frac{1}{2}}\\
T_{e}(t) & =\frac{\left\{ \sqrt{2t\left(C_{AR}^{(e)}+C_{AR}^{(h)}\right)+N_{e,0}^{-2}}\left[3k_{B}T_{e}N_{e,0}\left(C_{AR}^{(e)}+C_{AR}^{(h)}\right)+2E_{g}C_{AR}^{(e)}N_{e,0}\right]-2E_{g}C_{AR}^{(e)}\right\} }{3k_{B}\left(C_{AR}^{(e)}+C_{AR}^{(h)}\right)}
\end{cases}
\]
In that calculation, the band gap energy is taken constant. 

\begin{figure}[h]
\begin{centering}
\includegraphics[width=8cm,bb = 0 0 200 100, draft, type=eps]{home/thibault/Documents/LaAPT/Flaps2D/09-EnergyConservation/20140314-ValidationOfAugerCoupledTeNe.eps}
\par\end{centering}

\caption{\label{fig:ErrorOnAugerCoupled}Calculation of the error (\%) made
on the solutions $N_{e,h}(t)$ and $T_{e,h}(t)$ for $N_{e,h}(t=0)=10^{27}$
m$^{-3}$, $T_{e,h}(t=0)=80$ K. }
\end{figure}
Figure \ref{fig:ErrorOnAugerCoupled} shows the relative error $\zeta$
between the numerical and analytical results, calculated for $T_{e,h}$
and $N_{e,h}$ using the following formula
\[
\zeta(T_{e})=100\times\frac{T_{e}^{(num)}-T_{e}^{(exact)}}{T_{e}^{(exact)}}
\]


The first conclusion is that the Auger recombination is very well
solved. Calculating on a time range of $10^{-8}$ s with a timestep
of $dt\sim10^{-17}$ s, the relative error is lower than 0.1 \% on
the carrier temperature, and lower than $5\times10^{-3}$ \% on the
carrier density. The multiple time scale of the problem is taken into
account by changing the timestep during the calculations. In this
calculation, the timestep has been changed at $t=10^{-12}$ s, $t=10^{-11}$s,
and $t=10^{-10}$ s by multiplying $dt'\leftarrow10\times dt$. We
can directly observe the negligible consequences of the timestep modification
as a function of time. 

\begin{figure}[h]
\begin{centering}
\includegraphics[width=8cm,bb = 0 0 200 100, draft, type=eps]{home/thibault/Documents/LaAPT/Flaps2D/09-EnergyConservation/CheckAuger/20140312-Energy.eps}
\par\end{centering}

\caption{\label{fig:AugerEnergyConservation}Energy of the free-carriers as
a function of time after the maximum of the laser pulse. $F=10$ $J.m^{-2}$,
$\tau=40$ fs, $\lambda=343$ nm. It is worth to note that considering
Fermi-Dirac statistics leads to the same result. }
\end{figure}
Figure \ref{fig:AugerEnergyConservation} presents the conservation
of the carrier energy after their excitation by 1 photon absorption.
The Fermi-Dirac statistics are here taken into account. It is shown
that the free-carrier energy is well conserved. 


\section{\label{sec:Control-of-the}Control of the energy conservation}

Since the mesh is not regular and that finding analytical solution
to the equation using an arbitrary contour in a reasonable duration
is beyond the capabilities of the author, it has been decided to study
the energy conservation as a function of time. The timescale of the
processes investigated by this code is near from the nanosecond reaction
of the system, which is 6 orders of magnitude larger than the laser
pulse duration (of about 40 fs to 500 fs). The validity of the time-integrator
is thus of major importance in this multi-scale problem, the connexion
to the spatial steps is thus straightforward. 


\subsection{Energy of the carriers}

Since we model an external laser pulse irradiating the sample by its
intensity evolution with time and space through the function $I(t,x,y)$,
then we will express the total energy considering the invariance by
translation with the $z$ axis, normal to the orientation of this
document plane. 

The total energy study will be expressed in $J/m$. Actually, the
energy of the carriers in 2D is defined by 
\[
dE_{e}\left[J.m^{-3}\right]=dE_{e,kin}+dE_{pot}=d\left(C_{e}T_{e}\right)+d\left(E_{g}n_{e}\right)
\]
where $\mathcal{V}$ {[}m$^{2}${]} is the control volume of each
2D cell of the mesh, $E_{e,kin}$ expresses the thermal energy of
conduction band electrons, $E_{e,pot}$ expresses the potential energy
of the conduction band electrons. This latter expression can be also
written for holes. However, since the origin of the potential energies
is taken at the edge of the valence band, then holes have a potential
energy equal to 0. Thus we can write 
\[
dE_{h}\left[J.m^{-3}\right]=dE_{h,kin}+dE_{h,pot}=d\left(C_{h}T_{h}\right)
\]


Integrating over the two dimensions of the problem leads to write
the total energy $\xi_{e,h}$ at a given instant in the whole sample
volume by 
\begin{alignat*}{1}
\xi_{e}\,[J.m^{-1}] & =\iint dE_{e}\, d\mbox{\ensuremath{\mathcal{V}}}\\
 & =\iint\left(T_{e}dC_{e}+C_{e}dT_{e}+dE_{g}n_{e}+E_{g}dn_{e}\right)\, d\mathcal{V}
\end{alignat*}
 
\begin{alignat*}{1}
\xi_{h}\,[J.m^{-1}] & =\iint dE_{h}\, d\mbox{\ensuremath{\mathcal{V}}}\\
 & =\iint\left(T_{h}dC_{h}+C_{h}dT_{h}\right)\, d\mathcal{V}
\end{alignat*}
Band gap is a function of lattice temperature $T_{s}$ and carrier
density $n_{e}$ expressed via \cite{Driel1987} 
\[
E_{g}(T_{s},n_{e})=1.17-4.73\times10^{-4}\frac{T_{s}^{2}}{T_{s}+636}-1.5\times10^{-10}n_{e}^{1/3}
\]



\subsection{Control of the laser energy}

Now we have chosen the unity to express the energy in the problem,
we express now the total laser energy is by the expression
\[
\xi_{sample}[J.m^{-1}](y)=\iint I(t,x,y)\, dt\, dx
\]
integrated over the diameter of the laser spot. 

However, the shape of the surface is not flat, and the surface reflectivity
is taken into account by the analytical solutions of Mie scattering,
assuming a scattering on a 3D cylinder from which a plane section
is studied: the intersection between the study plane and the entire
3D space. Consequently, the calculation of the overall reflectivity
and transmittance by the problem shape is rather difficult. 


\subsubsection{Calculation of the tip absorptance}

We try to solve this problem by studying the energy absorbed by the
nanometric sized sample via 
\[
\xi_{absorbed}[J.m^{-1}]=\iiint\alpha I(t,x,y)\, dt\, dx\, dy
\]


In the conditions of the experiments ($\tau=40$ fs, $F=10$ J/m$^{2}$,
$\lambda=343$ nm), the order of magnitude of $\xi_{absorbed}$ is
$10^{-4}$ J. In order to calculate the overall absorbed laser energy,
we define 
\[
A=\frac{\xi_{absorbed}}{\xi_{laser}}
\]
where 
\[
\xi_{laser}=F\times\mathcal{V}
\]


\begin{table}[h]
\begin{centering}
\begin{tabular}{ccc|cc|c|c}
 &  & \multicolumn{1}{c}{} &  & \multicolumn{1}{c}{} & \multicolumn{1}{c}{} & \tabularnewline
\hline 
\hline 
{\small{}$M$} & {\small{}$N$} & {\small{}MeshShift} & {\small{}$\mathcal{V}$} & {\small{}$dt$} & {\small{}$\xi_{absorbed}$ (J)} & {\small{}$A(\%)$}\tabularnewline
\hline 
{\small{}11} & {\small{}11} & {\small{}1} & {\small{}$6.7004\times10^{-12}$} & {\small{}$10$ as} & {\small{}$5.2097\times10^{-11}$} & {\small{}$77.75$}\tabularnewline
{\small{}11} & {\small{}11} & {\small{}4} & {\small{}$6.72711\times10^{-12}$} & {\small{}``} & {\small{}$5.5004\times10^{-11}$} & {\small{}$81.76$}\tabularnewline
{\small{}51} & {\small{}51} & {\small{}4} & {\small{}$6.73268\times10^{-12}$} & {\small{}``} & {\small{}$6.1999\times10^{-11}$} & {\small{}$92.08$}\tabularnewline
{\small{}101} & {\small{}101} & {\small{}4} & {\small{}$6.73404\times10^{-12}$} & {\small{}``} & {\small{}$6.4153\times10^{-11}$} & {\small{}$95.26$}\tabularnewline
{\small{}101} & {\small{}101} & {\small{}10} & {\small{}$6.73415\times10^{-12}$} & {\small{}``} & {\small{}$6.4271\times10^{-11}$} & {\small{}$95.44$}\tabularnewline
{\small{}101} & {\small{}101} & {\small{}20} & {\small{}$6.73432\times10^{-12}$} & {\small{}``} & {\small{}$6.4122\times10^{-11}$} & {\small{}$95.21$}\tabularnewline
{\small{}151} & {\small{}151} & {\small{}4} & {\small{}$6.73439\times10^{-12}$} & {\small{}``} & {\small{}$6.4955\times10^{-11}$} & {\small{}$96.45$}\tabularnewline
{\small{}151} & {\small{}151} & {\small{}20} & {\small{}$6.7345\times10^{-12}$} & {\small{}``} & {\small{}$6.4528\times10^{-11}$} & {\small{}$95.81$}\tabularnewline
\hline 
\hline 
 &  & \multicolumn{1}{c}{} &  & \multicolumn{1}{c}{} & \multicolumn{1}{c}{} & \tabularnewline
\end{tabular}
\par\end{centering}

\caption{\label{tab:Optical-energy}Optical energy as a function of the mesh
precision for an homogeneous energy distribution $\nabla I(x,y)=0$,
where $F=10$ $J.m^{-2}$, and $\tau=40$ fs. }
\end{table}
Table \ref{tab:Optical-energy} presents the total optical energy
deposited in the sample as a function of the mesh properties in case
of an homogeneous distribution of laser energy. This allows to calculate
the minimal mesh and timestep parameters required to reach convergence
in the numerical sample volume and the convergence of the intensity
calculation. 


\subsubsection{Calculation of the incident laser energy}

The intensity at the tip apex interface is defined by the following
formula
\[
E_{laser}\left[J\right]=\iint\overrightarrow{S}\cdot\overrightarrow{n}\, d\mathcal{A}\, dt
\]
where $\overrightarrow{S}$ is the Poynting vector, $\overrightarrow{n}$
is the normal to the 3D surface of the cone, and $\mathcal{A}$ is
the area element defined around the normal vector. The intensity is
defined by the flux of the Poynting vector by the surface of area
$d\mathcal{A}$. In the plane of the simulation $\left(xOy\right)$,
we can write the following relation 
\[
\overrightarrow{S}\cdot\overrightarrow{n}=I\left[x,y_{\partial\Omega}\right]\cos\theta_{\partial\Omega}
\]
where 
\[
\tan\theta_{\partial\Omega}=\frac{\partial x}{\partial y}.
\]
Since we consider only two dimensions with an invariance by translation
along the coordinate $z$, we can build the quantity
\[
\xi_{laser}\left[J.m^{-1}\right]=\iint\overrightarrow{S}\cdot\overrightarrow{n}\, d\mathcal{L}\, dt
\]
using the contour of the cone in the plane of the model $\left(xOy\right)$,
where the projection of the Poynting vector at the interface is known. 

The arc length along the curve $x(y)$ is given by 
\[
\mathcal{L}=\int\left|\frac{\partial x}{\partial y}\right|\, dy
\]
Therefore, we can write
\[
\xi_{laser}\left[J.m^{-1}\right]=\iint I\left[x_{\partial\Omega}\left(y\right),y,t\right]\cos\left(\arctan\frac{\partial x}{\partial y}\right)\left|\frac{\partial x}{\partial y}\right|\, dy\, dt.
\]
For a constant intensity defined by $I(x,y,t)=I_{0}e^{-\frac{1}{2}\left(\frac{t}{\sigma_{t}}\right)^{2}}$,
$\sigma_{t}=\frac{\tau}{2\sqrt{2\ln2}}$, $I_{0}=\frac{F}{\tau}\sqrt{\frac{4\ln2}{\pi}}$,
$F=10$ $J.m^{-2}$, and $\tau=40$ fs, it is found 


\subsubsection{Verification of the calculated intensity}

\[
\xi_{laser}=\iint\left(\sigma_{1}+\alpha_{fcr,e}+\alpha_{fcr,h}\right)I_{0}(t,x,y)\, d\Omega\, dt
\]
\[
\xi_{opt}=\iint\left(\sigma_{1}+\alpha_{fcr,e}+\alpha_{fcr,h}\right)I(t,x,y)\, d\Omega\, dt
\]
We define the numerical trust value which corresponds to the ratio
between the energy absorbed contained the solid, with the energy given
by the laser intensity distribution. 
\[
\eta=\frac{\xi_{solid}}{\xi_{opt}}
\]
We also define the absorptivity by 
\[
\overline{A}=\frac{\xi_{opt}}{\xi_{laser}}
\]
 which is the effective absorptivity of the laser pulse considering
the shape of the cone. 


\paragraph{Results}

Table \ref{tab:EnergyBalanceMie} shows the results for 343 nm using
different meshes. It is clear that :
\begin{itemize}
\item The ``numerical trust index'' $\eta$ demonstrates that the size
of the mesh is of major importance for obtaining a good energy conservation
between the field and the solid. But why, since the equations are
using the same mesh ? It can be due to the fact that the strongest
absorption zone at 343 nm is at the boundary, and that the first energy
contained in the cells of the boundary is neglected. Thus, at 515
nm, we shall need less cells to reach the energy balance between the
fields and the solid. 
\item the effect of the meshgrid is negligible in the energy balance. It
will take its importance only for the energy diffusion conservation.
\item A problem is identified in the definition of the optical absorption
which simply cannot be so small. Indeed, at $\lambda=343$ nm, the
penetration depth of the field is nearly 9 nm, which is clearly smaller
than the size of the cone. Thus, it is not possible that the absorbed
energy is so small. Its definition might be wrong or reveals a numerical
problem. 
\end{itemize}
\begin{table*}[h]
\begin{centering}
\begin{tabular}{ccc|cc|ccc|ccc|cc}
 &  & \multicolumn{1}{c}{} &  & \multicolumn{1}{c}{} &  &  & \multicolumn{1}{c}{} &  &  & \multicolumn{1}{c}{} &  & \tabularnewline
\hline 
\hline 
{\scriptsize{}$M$} & {\scriptsize{}$N$} & {\scriptsize{}MeshShift} & {\scriptsize{}$\mathcal{V}$} & {\scriptsize{}$dt$} & {\scriptsize{}$\lambda$ (nm)} & {\scriptsize{}Pola} & {\scriptsize{}$\phi$ (rad)} & {\scriptsize{}$\xi_{laser}$ (mJ.m$^{-1}$)} & {\scriptsize{}$\xi_{opt}$ (mJ.m$^{-1}$)} & {\scriptsize{}$\xi_{solid}$ (mJ.m$^{-1}$)} & {\scriptsize{}Trust $\eta$ (\%)} & {\scriptsize{}$\overline{A}=1-R-T$ (\%)}\tabularnewline
\hline 
{\scriptsize{}101} & {\scriptsize{}101} & {\scriptsize{}4} & {\scriptsize{}$6.73404\times10^{-12}$} & {\scriptsize{}``} & {\scriptsize{}$343$} & {\scriptsize{}TM} & {\scriptsize{}$0$} & {\scriptsize{}$7.384$} & {\scriptsize{}$33.71\times10^{-3}$} & {\scriptsize{}$13.75\times10^{-3}$} & {\scriptsize{}$40.7$} & {\scriptsize{}$0.46$}\tabularnewline
{\scriptsize{}101} & {\scriptsize{}101} & {\scriptsize{}10} & {\scriptsize{}$6.73415\times10^{-12}$} & {\scriptsize{}``} &  &  &  &  &  &  &  & \tabularnewline
{\scriptsize{}101} & {\scriptsize{}101} & {\scriptsize{}20} & {\scriptsize{}$6.73432\times10^{-12}$} & {\scriptsize{}``} &  &  &  &  &  &  &  & \tabularnewline
{\scriptsize{}151} & {\scriptsize{}151} & {\scriptsize{}4} & {\scriptsize{}$6.73439\times10^{-12}$} & {\scriptsize{}``} & {\scriptsize{}$343$} & {\scriptsize{}TM} & {\scriptsize{}$0$} & {\scriptsize{}$7.384$} & {\scriptsize{}$29.53\times10^{-3}$} & {\scriptsize{}$18.07\times10^{-3}$} & {\scriptsize{}$61.2$} & {\scriptsize{}$0.40$}\tabularnewline
{\scriptsize{}251} & {\scriptsize{}251} & {\scriptsize{}4} & {\scriptsize{}$6.73462\times10^{-12}$} & {\scriptsize{}``} & {\scriptsize{}$343$} & {\scriptsize{}TM} & {\scriptsize{}$0$} & {\scriptsize{}$7.384$} & {\scriptsize{}$27.093\times10^{-3}$} & {\scriptsize{}$21.808\times10^{-3}$} & {\scriptsize{}$80.49$} & {\scriptsize{}$0.37$}\tabularnewline
{\scriptsize{}351} & {\scriptsize{}351} & {\scriptsize{}40} & {\scriptsize{}$6.73472\times10^{-12}$} & {\scriptsize{}``} & {\scriptsize{}$343$} & {\scriptsize{}TM} & {\scriptsize{}$0$} & {\scriptsize{}$7.384$} & {\scriptsize{}$26.340\times10^{-3}$} & {\scriptsize{}$23.061\times10^{-3}$} & {\scriptsize{}$88.95$ (no diff.)} & {\scriptsize{}$0.356$}\tabularnewline
\hline 
\hline 
 &  & \multicolumn{1}{c}{} &  & \multicolumn{1}{c}{} &  &  & \multicolumn{1}{c}{} &  &  & \multicolumn{1}{c}{} &  & \tabularnewline
\end{tabular}
\par\end{centering}

\caption{\label{tab:EnergyBalanceMie}Optical energy as a function of the mesh
precision for an inhomogeneous energy distribution $\nabla I(x,y)\protect\neq0$,
where $F=10$ $J.m^{-2}$, and $\tau=40$ fs. ``Trust'' is defined
by $\eta=\frac{\xi_{solid}}{\xi_{opt}}$.}
\end{table*}



\section{Generation and control of the mesh}


\paragraph{Definition of the contour}


\paragraph{Regularization of the mesh}

The mesh is built by defining a contour and solving 
\[
\Delta\left(\begin{array}{c}
x\\
y
\end{array}\right)=0
\]
 by a simple iterative method
\[
x_{i}=\frac{1}{4}\left(x_{i-1,j-1}+x_{i-1,j+1}+x_{i+1,j-1}+x_{i+1,j+1}\right)
\]


\[
y_{i}=\frac{1}{4}\left(y_{i-1,j-1}+y_{i-1,j+1}+y_{i+1,j-1}+y_{i+1,j+1}\right)
\]
Round about 70 000 iterations are required to reach convergence of
the mesh. 

This method is optimum on a single core since convergence rate of
this technique is rather good (how to quantify ?) in respect of the
global calculation time. 


\section{Optimization of the code}

Optimization is a big word for telling many things. First, we have
designed a model which is actually already simplified to decrease
the amount of necessary operations. In this section, purpose is to
perform optimizations over the code, especially over the way to perform
the operations on a single processor, and to decrease the waste of
operations. 


\subsection{Sequential code: checking the bottlenecks}

Using the Intel Loop Observer software (or GNU profiler), we can monitor
the number of calls to the loops and number of calls to the functions. 

By performing a first test during one night, I already figured out
that a function was using 25\% of the calculation time (sequential)
despite that it was not necessary. The number of operations was found
to be $10^{12}$ after one night. Actually, this operation was required
once, and then replaced by calls to memory. But is this call to memory
really saves time instead of performing such simple operation each
time ? 

In the next section, we focus on optimizing the way to perform calculations
with several cores, to parallelize the operations. First, we will
define the tools used to quantify the efficiency of a parallel code,
and we will also perform some verification allowing to avoid the unnecessary
operations. 


\subsection{Calculation of the code velocity}

\[
CPU=\frac{N_{timesteps}}{t_{current}-t_{initial}}\times N_{threads}
\]
where $t_{current}$ is the CPU time captured just before performing
the outputs, and $t_{initial}$ is the CPU time captured at the initialization
of the code. Hence, CPU is supposed to be constant within the execution
time of the code. 

The coefficient $CPU$ is built to be independent to number of cores
used. In this purpose, it becomes possible to calculate the speedup
$\eta$ in core-number-independent way. Of course, the speedup may
increase while increasing the number of cores, hopefully, until saturation. 

\begin{figure}[h]
\begin{centering}
\includegraphics[width=8cm,bb = 0 0 200 100, draft, type=eps]{home/thibault/Documents/Codes/Flaps/Optimization/20150208-SpeedUp.eps}
\par\end{centering}

\caption{CPU efficiency as a function of core number. Note that these calculations
were performed with a homogeneous energy source in the sample. }
\end{figure}


One understands that the definition of CPU is dependent to the frequency
of the CPU clock. Hence, it is not reasonable to compare the CPU quantity
between different computers. However, we can normalize the CPU by
the velocity of the code in sequential mode on a single computer in
order to build the speed up $\eta$, a quantity that it is intrinsic
to the efficiency of the code. 
\[
\eta(N)=\frac{CPU_{N-cores}}{CPU_{1-core}}
\]
Hence, it looks that speedup of our parallelization method is greater
than 1, but is lower than 2 with the available number of cores.


\subsection{Changing the order of the loops}

The main loop of the code is consuming 30\% of the calculation time
in sequential mode. By changing the order of the loops, and by parallelizing
the inner or outer ones, we can obtain serious improvement over the
speedup. Fig. \ref{fig:Optimization1} demonstrates a very strong
influence over the order of the loop iterations and distribution.
Despite that one solution looks to be more efficient than the others,
it is not sure if such optimization conserves the physical meaning
of the result. 

\begin{figure}[h]
\begin{centering}
\includegraphics[width=8cm,bb = 0 0 200 100, draft, type=eps]{home/thibault/Documents/Codes/Flaps/Optimization/20150208-OptimizeLoopOrder.eps}
\par\end{centering}

\caption{\label{fig:Optimization1}Qualitative increase of the efficiency as
a function of the main loop order in Flaps2D.}
\end{figure}



\subsection{Parallel code: reduce the number of conditional branches}

\noun{If} instruction is consuming a lot of time. Hence, it must be
removed, especially from loops. Problem is that a large code allows
to perform many types of calculations and to select various parameters.
How can we select sub-modes of calculations into a code without using
condition branches ? 

\begin{figure}[h]
\begin{centering}
\includegraphics[width=8cm,bb = 0 0 200 100, draft, type=eps]{home/thibault/Documents/Codes/Flaps/Optimization/20150208-SplitTheLoops.eps}
\par\end{centering}

\caption{\label{fig:Comparison-before-and}Comparison before and after reduction
of the IF branches inside parallelized loops. }
\end{figure}


Figure \ref{fig:Comparison-before-and} indicates that actually, cutting
the loops into one parallel loop per equation is not the best idea.
Indeed, the physical quantities are independent, since we use an explicit
solving scheme. Hence, increasing the number of parallel loops increase
the number of synchronization delays expected at the end of solving
each equation. It is then more efficient to parallelize one mesh node
per loop and to solve everything in this mesh node using the same
core. 


\subsection{Efficiency of the parallelization method}


\paragraph{Distribution of the loops on CPU cores}

The method of distributing the loops over the various cores has a
limited efficiency. However, we can try to optimize the way they are
distributed with the command \noun{omp\_schedule. }

\begin{figure}[h]
\begin{centering}
\includegraphics[width=8cm,bb = 0 0 200 100, draft, type=eps]{home/thibault/Documents/Codes/Flaps/Optimization/20150208-OptimizeLoopDistrib.eps}
\par\end{centering}

\caption{\label{fig:CPU-efficiency-as}CPU efficiency as a function of the
method to distribute the loops with OpenMP using 8 cores. }
\end{figure}
Figure \ref{fig:CPU-efficiency-as} presents the CPU efficiency evolution
as a function of the strategy to distribute the loops over the cores.
These calculations were performed only using 8 cores. 


\paragraph{Explicit integration VS implicit integration}

\begin{figure}[h]
\begin{centering}
\includegraphics[width=8cm,bb = 0 0 200 100, draft, type=eps]{home/thibault/Documents/Codes/Flaps/Optimization/20150208-ComparisonExplImpl.eps}
\par\end{centering}

\caption{\label{fig:ComparisonExplImpl}CPU efficiency as a function of the
method to distribute the loops with OpenMP using 8 cores. }
\end{figure}
Overall the optimization performed until now, Figure \ref{fig:ComparisonExplImpl}
presents the efficiency of the code for the implicit version and the
explicit version, including recent optimizations. The implicit version
of the code consumes round about 7 GB of RAM, and the explicit version
consumes \textasciitilde{} 400 MB. 


\paragraph{\emph{Which parallelization method can be investigated to strongly
increase the speedup ? }}
\begin{enumerate}
\item Go to GPU. 
\item Find OpenMP bottlenecks using OpenMP profiler program. 
\item Try the code on 16 cores CPU to verify scalability. 
\end{enumerate}
\bibliographystyle{unsrtnat}
\bibliography{/media/thibault/PHD/Travail/These/Articles/bibliographie_lue}

\end{document}
